% !PRINTED 599
\weberpage{599}
%<<CONT
degree $h_r$ of the intersection of $\Phi'$ and $\Psi_r$. The degree of $\Phi_r$
is therefore equal to $f g n' : h_r$.

If we now denote by $\varphi_r$ all the elements of $\Phi_r$, then all
the prime ideals $\mathfrak{p}\,|\,\varphi_r$ in $\Omega$ lead to the same prime ideal
$\mathfrak{p}_r'$ in $\Omega'$. But it remains to be shown that two different ones
of these systems $\Phi_1, \Phi_2$ lead to \gesperrt{different} prime ideals $\mathfrak{p}_1', \mathfrak{p}_2'$.

First of all it is evident that all the $\mathfrak{p}\,|\,\varphi_2$ are different from
the $\mathfrak{p}\,|\,\varphi_1$. For if $\mathfrak{p}\,|\,\varphi_1 = \mathfrak{p}\,|\,\varphi_2$, then also
$\mathfrak{p} = \mathfrak{p}\,|\,\varphi_2\,\varphi_1^{-1}$, so that $\varphi_2\,\varphi_1^{-1}$ is contained in $\boldsymbol{\Psi}$, and consequently
$\varphi_2$ in $\boldsymbol{\Psi}\varphi_1$, that is, in $\Phi_1$, contrary to the assumption. We
may therefore assume an integer $\alpha$ in $\Omega$ which is divisible by
$\mathfrak{p}\,|\,\varphi_1$, but by none of the ideals $\mathfrak{p}\,|\,\varphi_2$. Then none of the
numbers $\alpha\,|\,\varphi'$ is divisible by $\mathfrak{p}\,|\,\varphi_2$ either, since otherwise
$\mathfrak{p}\,|\,\varphi_1\,\varphi' = \mathfrak{p}\,|\,\varphi_2$, and hence, contrary to the assumption, $\varphi_2$ would be
contained in $\Phi_1$. The product $\alpha'$ of all the mutually different
$\alpha\,|\,\varphi'$, which is a number contained in $\Omega'$, is divisible by
$\mathfrak{p}_1'$, but not by $\mathfrak{p}_2'$. Consequently $\mathfrak{p}_1'$ is different from $\mathfrak{p}_2'$,
and it thus results that $e'$ and
no more \gesperrt{different} prime ideals $\mathfrak{p}'$ divide $\mathfrak{P}$.
We therefore set

\weq{3}{\mathfrak{P} = \mathfrak{p}_1'^{\,a_1}\,\mathfrak{p}_2'^{\,a_2}\cdots\mathfrak{p}_{e'}'^{\,a_{e'}},}

in which the exponents $a_1, a_2, \ldots, a_{e'}$ are natural numbers still
to be determined more closely.

Now, in order to decompose the prime ideals $\mathfrak{p}_r'$ in the field $\Omega$,
we can use the results of \S.\,160, 162 by simply letting $\Omega'$
take the place of $R$. If then $X_r'$ denotes the intersection
of $\Phi'$ with $X_r = \varphi_r^{-1}\,X\,\varphi_r$, then $X_r$ is the totality of all % [sic] source omits the prime; the next sentence correctly refers to $X_r'$
% [sic: The source says "Körpers" although $\Phi'$ denotes the group, not the field.]
substitutions of the field $\Phi'$ which satisfy the condition

\[
\tau\,|\,\chi_r' \equiv \tau \ (\mathrm{mod}\ \mathfrak{p}_r)
\]

If we therefore denote the degree of $X_r'$ by $g_r$, then
$\mathfrak{p}_r'$ is divisible by $\mathfrak{p}_r^{g_r}$, but by no higher power of $\mathfrak{p}_r$.

The intersection $\Psi_r'$ of $\Psi_r$ with $\Phi'$, whose degree we have
already denoted above by $h_r$, is the totality of all substitutions
$\psi_r'$ in $\Phi'$ which satisfy the condition

\[
\mathfrak{p}_r\,|\,\psi_r' = \mathfrak{p}_r
\]

and if we therefore decompose $\Phi'$ into the subsidiary groups

\weq{4}{\Phi' = \Psi_r'\varphi_{r,1}' + \Psi_r'\varphi_{r,2}' + \cdots + \Psi_r'\varphi_{r,e_r}} % [sic] source omits the prime on the final $\varphi$, unlike the two preceding terms

then

\weq{5}{n' = e_r\,h_r.}
